\section{Rest energy in the closure}
\label{sec:restclosure}

The closure of \cref{sec:closure} is the one dynamics that the arena closes by itself: the rate $\kappa$ of the shift operators
$\alpha_\pm=\partial_\tau\pm\kappa G$ becomes the angular velocity of a rotor. This section asks what that dynamics does to the energies of a
field point, the undivided energy $\tfrac12(|\bB|^2+|\bE|^2)$ and the rest energy, and whether it relates the rest energy to the rate. Two fields
are kept apart throughout: the \emph{momentum field} $\pi=\eta^{-1}p$, which is $m\alpha_+X$ in the clamped split-form interior, and the
\emph{frame field} $X$. The answers, in brief. The closure freezes the undivided energy, the rest energy and the Poynting vector of the momentum
field, for every Hamiltonian; it changes those of the frame field only by an oscillation or an unbounded drift, and selects no value. The rest
energy of the momentum field is exactly the amplitude of the potential of the duality angle, and the conserved total splits into a rate part and a
Lorentz-invariant frame part that exchange that amplitude. The constant rate of the interior is the limit of a heavy rotor. The rest flow of Dirac
matter and the binding rotation of the frame are two rotations at rest that no operator of the arena relates, so $mc^2=\hbar\kappa$ is an added
identification, and nothing in the algebra fixes the length $\ell$ in $\kappa=c/\ell$. A quantised rotor gives a rate ladder only under stated
assumptions.

\paragraph{Notation.} As in \cref{sec:closure}, $V=\R^{3,3}$, $G$ is the binding root, $G^2=-\id$, $G^{\!\top}=-G$, $\eta G=-G\eta$, a covector
$p$ is a column, and $z(Y)=\bB+i\bE$ for $Y\in V$ with components $(\bB;\bE)$ along $(e_1,e_2,e_3;e_4,e_5,e_6)$, so that $G$ acts as $i$
(\cref{thm:emfield}). With $q(z)=z\cdot z=|\bB|^2-|\bE|^2+2i\,\bB\cdot\bE$ put
\begin{equation}
\label{eq:restdefs}
E_6(Y):=\tfrac12|z(Y)|^2,\qquad E_0(Y):=\tfrac12\bigl|q\bigl(z(Y)\bigr)\bigr|,\qquad \mathbf S(Y):=\tfrac i2\,\overline{z(Y)}\times z(Y)=\bE\times\bB ,
\end{equation}
the undivided energy, the rest energy (the rest-frame energy density of \cref{sec:orth}) and the Poynting vector of $Y$. Since $\Btext=-\bB$,
$\mathbf S$ is minus the textbook Poynting vector; only $|\mathbf S|$ enters below. For the momentum field, $E_6(\pi)=\tfrac12|p|^2$ because $\eta$
is orthogonal, and $E_0(\pi)=\tfrac12|w(p)|$ with $w$ as in \eqref{eq:wdef}. The Lorentz-class inertia of \cref{thm:inertiareq}(ii) is written
$N=(\nu_1\eta^{-1}+\nu_2G\eta^{-1})/m$, $\nu:=\nu_1+i\nu_2$, its coefficients being the $(a,b)$ of that theorem. Units are $\hbar=c=1$ unless a
symbol is shown, $\tau$ is a time, and $\mathrm K(\mu)$, $\mathrm E(\mu)$ are the complete elliptic integrals of \cref{ex:drift}.

\subsection{What the closure conserves}
\label{sec:restconserve}

On $\mathfrak s^*$, the dual of the unique closure $\mathfrak s=\R\hat g\ltimes_GV$ of \cref{thm:closureunique}, with coordinates $(h,p)$,
$F_h=\partial F/\partial h$ and $\nabla F=\partial F/\partial p$, the Lie--Poisson bracket and the equations of an arbitrary Hamiltonian $H$ are
\begin{equation}
\label{eq:restLPgen}
\{F_1,F_2\}=p^{\!\top}\bigl(F_{2,h}\,G\nabla F_1-F_{1,h}\,G\nabla F_2\bigr),\qquad
\dot p=H_h\,G^{\!\top}p,\qquad \dot h=-\,p^{\!\top}G\,\nabla H .
\end{equation}
For $H=h^2/(2I_0)+\tfrac12p^{\!\top}Np$ these are the equations \eqref{eq:LProtor} of \cref{prop:LP}.

\begin{theorem}[The momentum field is frozen up to its duality phase]
\label{cl:frozen}
\begin{enumerate}[label=(\roman*),leftmargin=2.2em]
\item A function $f(p)$ is a Casimir of $\mathfrak s^*$ if and only if it is invariant under $p\mapsto e^{tG}p$, that is under the duality rotations
$z(\pi)\mapsto e^{-it}z(\pi)$ of the momentum field. Hence $E_6(\pi)$, $E_0(\pi)$ and the vector $\mathbf S(\pi)$ are Casimirs, while $\Re w$ and
$\Im w$ are not.
\item At $p\neq0$ the symplectic leaves are two-dimensional, so there are five independent local Casimirs. The real and imaginary parts of the
products $\bar z_jz_k$, $z=z(\pi)$, provide five, and $E_6(\pi)$, $E_0(\pi)$, $\mathbf S(\pi)$ account for four of them, subject to
$E_6(\pi)^2=E_0(\pi)^2+|\mathbf S(\pi)|^2$.
\item Along the flow of every Hamiltonian, $p(\tau)=e^{-\theta(\tau)G}p_0$ with $\dot\theta=H_h$ and $\theta(0)=0$; so
$z(\pi(\tau))=e^{i\theta(\tau)}z(\pi_0)$ and $w(p(\tau))=e^{2i\theta(\tau)}w(p_0)$.
\end{enumerate}
\end{theorem}

\begin{proof}
(i) For $f=f(p)$, \eqref{eq:restLPgen} gives $\{f,F\}=F_h\,p^{\!\top}G\nabla f=-F_h\,(Gp)\cdot\nabla f$, which vanishes for every $F$ exactly when $f$
is constant along $t\mapsto e^{tG}p$. Since $\eta G=-G\eta$, $\eta^{-1}e^{tG}p=e^{-tG}\pi$, on which $G$ acts as $i$. The quantities $|z|^2$,
$|z\cdot z|$ and $\bar z\times z$ are invariant under a phase, while $w=z\cdot z$ is multiplied by $e^{-2it}$. (ii) From \eqref{eq:restLPgen},
$\{p_i,p_j\}=0$ and $\{h,p_j\}=(Gp)_j$, so the Poisson tensor has rank two at $p\neq0$ and the leaves have codimension five in the
seven-dimensional $\mathfrak s^*$. The map $z\mapsto z\bar z^{\!\top}$ from $\C^3\setminus\{0\}$ onto the rank-one positive Hermitian matrices has the
phase circles as its fibres, so its image is five-dimensional and its nine real components contain five independent functions, all phase
invariant. The relation is the identity $u^2-|\mathbf S|^2=\tfrac14|q|^2$ preceding \eqref{eq:EBbound}. (iii) $\dot p=-H_hGp$ integrates to
$p=e^{-\theta G}p_0$, so $\pi=e^{\theta G}\pi_0$, and $w$ is quadratic in $z$.
\end{proof}

The theorem is fixed by the algebra. In substance it is a corollary of \cref{prop:LP}, whose coadjoint orbits are the cylinders
$\R\times\{e^{\theta G}p_0\}$; item (iii) is the phase law of \cref{thm:pendulum}, extended to every Hamiltonian, and the relation in (ii) is
inherited. What it adds is that the invariants of the momentum field are exactly the Casimirs, and their count. Already in the clamped interior
$p=e^{-\kappa\tau G}p_0$, so these functions of the momentum field were constant before the rotor was freed.

\begin{corollary}[What the closure can and cannot change]
\label{cl:nogen}
\begin{enumerate}[label=(\roman*),leftmargin=2.2em]
\item For every Hamiltonian on $\mathfrak s^*$, quadratic or not, of Lorentz class or not, $E_6(\pi)$, $E_0(\pi)$ and $\mathbf S(\pi)$ are constant
along every solution: a null momentum field stays null, and a momentum field with rest energy $E_0(\pi_0)$ keeps it.
\item The same functions of the frame field are not constants of motion. In the clamped split-form interior, where
$z(X)=c_-e^{-i\kappa\tau}+c_+e^{i\kappa\tau}$ (\cref{prop:orbitsEM}(ii)),
\[
\begin{gathered}
E_6(X)=\tfrac12\bigl(|c_-|^2+|c_+|^2\bigr)+\Re\bigl(\bar c_-\!\cdot c_+\,e^{2i\kappa\tau}\bigr),\\
q\bigl(z(X)\bigr)=c_-\!\cdot c_-\,e^{-2i\kappa\tau}+2\,c_-\!\cdot c_++c_+\!\cdot c_+\,e^{2i\kappa\tau},
\end{gathered}
\]
so $E_6(X)$ is constant if and only if $\bar c_-\!\cdot c_+=0$, and $c_-=(1,i,0)$, $c_+=(1,0,0)$ give $|q(z(X))|^2=9-8\sin^2\kappa\tau$. Under the
free-rotor drift of \cref{ex:drift} with $x(0)=0$, the frame field $X=e^{-\theta G}x$ has $E_0(X)=0$ at $\tau=0$ and
$E_0(X)=E_6(X)=\tfrac12k^2|\Delta x|^2$ after $k$ full turns, $\Delta x$ the displacement per turn of that example; after three turns this is
$\approx1.807$. On every rotating solution with Lorentz-class inertia and $R_0>0$, in the notation of \cref{thm:pendulum}, $E_6(X)$ and
$E_0(X)$ are unbounded.
\end{enumerate}
\end{corollary}

\begin{proof}
(i) Casimirs Poisson-commute with every Hamiltonian (\cref{cl:frozen}(i)). (ii) The two formulas are the expansions of $\bar z\cdot z$ and $z\cdot z$.
In the example $c_-\cdot c_-=0$ and $c_-\cdot c_+=c_+\cdot c_+=1$, so $|q|^2=|2+e^{2i\kappa\tau}|^2=5+4\cos2\kappa\tau$. Since $G$ acts as $i$,
$e^{-\theta G}$ multiplies $z$ by a phase, so $E_6(X)=E_6(x)$ and $E_0(X)=E_0(x)$. In \cref{ex:drift} the motion of $(\theta,h)$ and the velocity
$\dot x=e^{2\theta G}Np_0$ are periodic in $\theta$, so at $\theta=2\pi k$ the frame is $x=k\,\Delta x$, a real multiple of $e_1$, and
$q(z(x))=k^2|\Delta x|^2$. In general, by \cref{thm:freedrift} the frame advances by $\Delta x=\lambda e^{-\varphi_0G}Np_0$, $\lambda>0$, in each
period of the rate, so $x=x_0+n\Delta x$ after $n$ periods. As $z(Np_0)=\nu\,z(\pi_0)/m$, $q(z(\Delta x))=\lambda^2e^{-2i\varphi_0}\nu^2w(p_0)/m^2$,
which is nonzero because $|\nu w(p_0)|=R_0>0$; hence $|q(z(x))|$ and $|x|^2$ grow like $n^2$.
\end{proof}

Item (i) is an obstruction, under the hypothesis that the dynamics is a Hamiltonian flow on $\mathfrak s^*$: no such flow produces, removes or
changes the rest energy of the momentum field. Item (ii) is fixed, from \cref{prop:orbitsEM,thm:freedrift}. The first verdict therefore cuts both
ways. For the momentum field, $E_6(\pi)$ and its split $E_6(\pi)^2=E_0(\pi)^2+|\mathbf S(\pi)|^2$ label the leaf: the rest energy is carried by the
dynamics, not produced by it. For the frame field, to which the field-strength reading of $E_6$ and $E_0$ refers, the closure does change the rest
energy, by an oscillation or by an unbounded drift, but not by any mechanism that selects a value. Nor is $E_6(X)$ a function of
$E_6(\pi)$: on the orbits without momentum of \cref{prop:orbitsEM}(i) with $X\neq0$, $|p|^2=0<E_6(X)$. What moves is the duality angle $\theta$
and the rate $h/I_0$.

\subsection{The duality pendulum}
\label{sec:restpendulum}

From now on $N$ is of Lorentz class, $\nu w(p_0)=R_0e^{i\varphi_0}$ with $R_0\ge0$, and $\psi:=2\theta+\varphi_0$. By \cref{thm:inertiareq}(ii) the
frame energy is $\tfrac12p^{\!\top}Np=\Re(\nu w(p))/(2m)$, and by \cref{thm:pendulum} the rotor is the pendulum
\[
\frac{h^2}{2I_0}+U(\theta),\qquad U(\theta)=\frac{R_0}{2m}\cos\psi,\qquad I_0\ddot\theta=\frac{R_0}m\sin\psi,\qquad
R_0=|\nu|\,|w(p_0)|=2|\nu|\,E_0(\pi_0).
\]
The amplitude of the potential of the duality angle is therefore the rest energy of the momentum field, times $|\nu|/m$; its undivided energy
$E_6(\pi)$ and its Poynting vector $\mathbf S(\pi)$ do not enter.

\begin{proposition}[Small oscillations]
\label{cl:omega}
The stable equilibrium is $\psi=\pi$, where $U=-R_0/(2m)$. With $\psi=\pi+\xi$, $\ddot\xi=-\Omega^2\sin\xi$ and
\[
\Omega^2=\frac{2R_0}{mI_0}=\frac{4|\nu|\,E_0(\pi_0)}{mI_0} ;
\]
$\theta$ oscillates at the same frequency with half the amplitude of $\xi$, and at amplitude $0<\xi_0<\pi$ the period is
$4\,\mathrm K\bigl(\sin^2(\xi_0/2)\bigr)/\Omega$.
\end{proposition}

\begin{proof}
$\ddot\psi=2\ddot\theta=(2R_0/mI_0)\sin\psi$ and $\sin(\pi+\xi)=-\sin\xi$. The period follows from $\tfrac12\dot\xi^2=\Omega^2(\cos\xi-\cos\xi_0)$ with
the substitution $\sin(\xi/2)=\sin(\xi_0/2)\sin t$.
\end{proof}

The proposition is fixed, as a corollary of \cref{thm:pendulum}, which it linearises; the period is the inherited one of the pendulum
\cite{DLMF}. For $R_0=0$, \cref{cor:constrate} says when the rate is constant: with Lorentz-class inertia it is constant and nonzero exactly
when the momentum field is null, or $\nu=0$, and $h(0)\neq0$; then $\theta=h\tau/I_0$ is an exact flat direction. For duality-class inertia, $[N,G]=0$, the rate is constant for all data
(\cref{thm:inertiareq}(iii), \cref{thm:drift}), and the potential $U(\theta)=\tfrac12p_0^{\!\top}e^{\theta G}Ne^{-\theta G}p_0=\tfrac12p_0^{\!\top}Np_0$ is
constant, in general not zero. The remaining case is the following.

\begin{corollary}[A constant rate with a rest energy]
\label{cl:flat}
Let $N$ be of Lorentz class with $\nu\neq0$, and let $E_0(\pi_0)>0$. The rate $h/I_0$ is constant along a solution if and only if $h(0)=0$ and
$\psi(0)\in\pi\Z$. There the rate vanishes and the frame moves at the constant velocity $Np_0$, with $|Np_0|=|\nu|\,|p_0|/m\neq0$.
\end{corollary}

\begin{proof}
$\dot h=(R_0/m)\sin\psi$ with $R_0>0$. If $h$ is constant, $\sin\psi\equiv0$, so $\theta$ is constant and $h=I_0\dot\theta=0$; conversely these points are
equilibria. With $\theta\equiv0$, \cref{thm:drift} ($N_c=0$) gives $\dot x=Np_0$, and $N^2=|\nu|^2\id/m^2$ because $\eta^{-1}G\eta^{-1}=-G$.
\end{proof}

\begin{remark}[A pseudo-Goldstone reading, and its limits]
\label{cl:goldstone}
On a leaf, $\theta$ is the angle of the duality group $U(1)_G$ of the momentum field, which a nonzero rest energy breaks to $\{\pm\id\}$ on the
invariants (\cref{sec:orth}). In the duality class it is a symmetry of $H$; in the Lorentz class it is broken on each leaf in proportion to
$E_0(\pi)$, and exactly
\[
I_0\,\Omega^2=U''(\theta_{\min})=\frac{2R_0}m=\frac{4|\nu|\,E_0(\pi_0)}m .
\]
The first equality defines $\Omega$ for every oscillator with one degree of freedom; the content specific to the closure is the second, which is
\cref{thm:pendulum}. With $I_0$ in the place of $f_\pi^2$ and $|\nu|E_0(\pi_0)/m$ in the place of the explicit breaking, the relation has the
shape of the Gell-Mann--Oakes--Renner relation for a pseudo-Goldstone boson \cite{GellMannOakesRenner68,Weinberg72}, as every pendulum does.
That is a reading, and its limits are definite. The breaking is explicit, by the inertia, not spontaneous: there is no vacuum manifold, and
$E_0(\pi)$ is a label fixed by the data (\cref{cl:nogen}(i)), not an order parameter; and by the bound \eqref{eq:EBbound} the field-point
invariants have no degenerate family of lowest-energy configurations among which to choose. With one degree of freedom, ``mass'' here means a frequency. What survives the reading is exact: on a null
momentum field the angle is a zero mode, and a rest energy of the momentum field gives it the frequency $\Omega\propto\sqrt{E_0(\pi_0)}$.
\end{remark}

\subsection{The energy budget: a rate part and a frame part}
\label{sec:restbudget}

\begin{proposition}[Rate part and frame part]
\label{cl:budget}
With Lorentz-class inertia, along every solution:
\begin{enumerate}[label=(\roman*),leftmargin=2.2em]
\item $H=K_{\rm rot}+U$ with $K_{\rm rot}=\tfrac12I_0\kappa^2\ge0$, $\kappa=h/I_0$, $U=\tfrac{R_0}{2m}\cos\psi=\Re(\nu\,w(p))/(2m)$, and
$\dot K_{\rm rot}=-\dot U=\tfrac{R_0}m\,\kappa\sin\psi$;
\item in the rotating regime $K_{\rm rot}$ runs through exactly $\bigl[H-\tfrac{R_0}{2m},\,H+\tfrac{R_0}{2m}\bigr]$ in every period of the rate, so the
energy exchanged between the two parts is $R_0/m=2|\nu|E_0(\pi_0)/m$; in the librating regime it runs through $\bigl[0,\,H+\tfrac{R_0}{2m}\bigr]$;
\item $K_{\rm rot}$ and $U$ are separately invariant under $\lG$ acting by $(h,p)\mapsto(h,\Lambda^{-\top}p)$, $\Lambda=e^r$, $r\in\lG$, and the flow
commutes with this action;
\item $(\theta,h)(\tau)$ depends on $p_0$ only through $\nu w(p_0)$; $E_6(\pi_0)$ and $\mathbf S(\pi_0)$ do not enter the motion.
\end{enumerate}
\end{proposition}

\begin{proof}
(i) $\dot K_{\rm rot}=h\dot h/I_0$ and $\dot U=-(R_0/2m)\,2\kappa\sin\psi$, with $\dot h=(R_0/m)\sin\psi$. (ii) In rotation $\dot\psi=2\kappa$ never
vanishes and $\cos\psi$ covers $[-1,1]$ in each period; in libration $K_{\rm rot}=0$ at the turning points and is largest at $\psi=\pi$. (iii) The
Lorentz requirement $rN+Nr^{\!\top}=0$ of \cref{def:inertiareq} integrates to $\Lambda^{-1}N\Lambda^{-\top}=N$, so $U(\Lambda^{-\top}p)=U(p)$; since
$[r,G]=0$, $\Lambda^{-\top}$ commutes with $G^{\!\top}$, and both equations \eqref{eq:LProtor} are preserved. (iv) $\eta^{-1}\Lambda^{-\top}=\Lambda\eta^{-1}$
because $\Lambda\in O(3,3)$, and $\lG\cong\so{3,\C}$ acts on $z$ complex-linearly and preserves $q$ (\cref{thm:GLorentz}), so
$w(\Lambda^{-\top}p)=w(p)$; the equations for $(\theta,h)$ contain only $R_0$ and $\varphi_0$.
\end{proof}

Item (i) restates the energy $E_\theta=H$ of \cref{thm:pendulum}; items (ii)--(iv) are fixed by the algebra. The conserved total thus
differentiates into two Lorentz-invariant parts, a rate part and a frame part, whose exchange is fixed by the rest energy of the momentum field
alone; neither part is the undivided energy $E_6(\pi)$, which is inert.

\paragraph{Units, and the factor one half.} Take $[\tau]=\mathsf T$ and $[\kappa]=\mathsf T^{-1}$, as in \eqref{eq:kappa}, and $[X]=\mathsf L$, a
convention by \cref{cl:scalings}(i) under which $m$ is a mass. Then $[I_0]=\mathsf{ML}^2$, $[h]=\mathsf{ML}^2\mathsf T^{-1}$ is an action,
$[p]=\mathsf{MLT}^{-1}$, $[N]=\mathsf M^{-1}$ and $[H]=[K_{\rm rot}]=[U]=\mathsf{ML}^2\mathsf T^{-2}$; $E_0(\pi)$ is a squared momentum and
$|\nu|E_0(\pi)/m$ an energy. With the clamped value $\kappa=c/\ell$ and $M_{\rm rot}:=I_0/\ell^2$, a mass,
\[
K_{\rm rot}=\tfrac12I_0\kappa^2=\tfrac12M_{\rm rot}c^2 .
\]
Reading this as the ``$c^2$ part'' of the conserved total is a reading: it is literally the Newtonian kinetic energy of a mass $M_{\rm rot}$ on a
circle of radius $\ell$ at rim speed $c$, outside the domain of Newtonian kinetics. It is not $E=mc^2$. The factor $\tfrac12$ comes from an energy
quadratic in the rate, where Planck's relation is linear; $M_{\rm rot}$ is not tied to $m$, since the rescaling $(I_0,m)\mapsto(sI_0,tm)$ recalled in \cref{sec:restscale} changes the two
independently; and $\kappa=c/\ell$ holds only
for the clamped rotor.

\subsection{When the rate is constant: the multiscale criterion}
\label{sec:restmultiscale}

\Cref{cor:multiscale} answers when the rate is nearly constant. In terms of the two parts of the energy its parameter is
\[
\varepsilon=\frac{R_0}{mI_0\kappa_0^2}=\frac{U_{\max}}{K_{\rm rot}(0)},\qquad U_{\max}=\frac{R_0}{2m}=\frac{|\nu|\,E_0(\pi_0)}m ,
\]
the ratio of the rest-energy amplitude of the momentum field to the initial rotor energy. If $c$ is the fixed conversion of \cref{thm:onec}, a
varying rate $\kappa=c/\ell$ is a varying length $\ell(\tau)=c/\kappa(\tau)$, and the law $\kappa^2=\kappa_0^2[1+\varepsilon(\cos\psi_0-\cos\psi)]$ of
\cref{cor:multiscale} gives it the relative spread $\varepsilon+O(\varepsilon^2)$ along a rotating solution. With $\kappa_0=c/\ell$, the multiscale
requirement \eqref{eq:multiscale} becomes
\[
\tfrac12M_{\rm rot}c^2\ \gg\ \frac{|\nu|\,E_0(\pi_0)}m :
\]
for ``$c/\ell$'' to be constant, the rotor energy must dominate the rest-energy amplitude of the momentum field. The clamped split-form interior
of \cref{sec:neutral} is the limit $\varepsilon\to0$ (\cref{prop:groupred,cor:multiscale}); the positive-form interior of \cref{sec:shifted} has
duality-class inertia and an exactly constant rate (\cref{thm:inertiareq}(iii)). Read in this way, the ``symmetric'' $c^2/\ell^2=(\kappa G)^2$ of
\cref{thm:EAdichotomy} is constant only insofar as the differentiated part is negligible beside the rotor; $\varepsilon$ is an input, and the
reading is ours.

\subsection{Two rotations at rest}
\label{sec:resttworot}

\begin{proposition}[Two rotations at rest]
\label{cl:tworot}
\begin{enumerate}[label=(\roman*),leftmargin=2.2em]
\item For the first-order real form $\dot\Psi=-JH\Psi$ of free Dirac evolution of \cref{sec:dirac} at $\bp=0$, $H(0)=me_1$ by \eqref{eq:Hp} and $Je_1=T=e_1e_2e_3$ with $T^2=-\id$,
so $\Psi(\lambda)=\cos(m\lambda)\Psi_0-\sin(m\lambda)\,T\Psi_0=e^{-m\lambda T}\Psi_0$: a rotation of angular frequency $m$, that is $mc^2/\hbar$.
\item In the split-form interior the orbits $X(\tau)=e^{-\kappa G\tau}X_0$ satisfy $\alpha_+X=0$ and $\ddot X=-\kappa^2X$: rotations of angular
frequency $\kappa=c/\ell$, with $z\mapsto e^{-i\kappa\tau}z$ and $q\mapsto e^{-2i\kappa\tau}q$. The orbits $X(\tau)=e^{\kappa G\tau}X_0$ are duality
rotations of the other sense at the same rate, and carry the momentum $\alpha_+X=2\kappa GX\neq0$.
\end{enumerate}
\end{proposition}

\begin{proof}
(i) $\dot\Psi=-JH(0)\Psi$ and $JH(0)=mJe_1=mT$, since $e_1$ commutes with $J=e_2e_3$. (ii) $\tfrac{d}{d\tau}e^{\mp\kappa G\tau}=\mp\kappa Ge^{\mp\kappa G\tau}$
and $G^2=-\id$.
\end{proof}

Item (i) is inherited: it is the realification of $\psi=e^{-imc^2t/\hbar}\psi_0$, de Broglie's internal clock \cite{deBroglie1925}. Item (ii)
restates \cref{prop:orbitsEM}(i)--(ii). Duality acts along both senses of the frame's rotation at the same rate, so the sense without momentum
is singled out by the momentum, not by duality.

\begin{proposition}[No operator relates the two rates]
\label{cl:norel}
\begin{enumerate}[label=(\roman*),leftmargin=2.2em]
\item $\Ad_{e^{\vartheta T}}$ fixes $e_1,e_2,e_3$ and maps a spatial generator $e_a$ to $\cos2\vartheta\,e_a+\sin2\vartheta\,Te_a$, which is a vector
only for $\vartheta\in\tfrac\pi2\Z$. The Dirac rest flow is therefore not a motion of $V$, and its generator $T$, of grade three, is not in $\so{3,3}$.
\item $(e^{\theta G})^{\!\top}\eta\,e^{\theta G}=\cos2\theta\,\eta+\sin2\theta\,\eta G$, so $e^{\theta G}$ lies in $O(3,3)$, and lifts to
$\mathrm{Pin}(3,3)$, only for $\theta\in\pi\Z$.
\end{enumerate}
Hence neither one-parameter group is the image of the other under the spin homomorphism, and this route supplies no relation between $m$
and $\kappa$.
\end{proposition}

\begin{proof}
(i) $T$ commutes with the temporal and anticommutes with the spatial generators (\cref{lem:blade}), so $e^{\vartheta T}e_ae^{-\vartheta T}=e^{2\vartheta T}e_a$,
and $Te_a$ has grade four. (ii) Expand with $e^{\theta G}=\cos\theta+\sin\theta\,G$, using that $G^{\!\top}\eta=\eta G$ is symmetric and
$G^{\!\top}\eta G=-\eta$.
\end{proof}

The proposition is an obstruction to one route only, a map of the arena carrying one rate to the other; item (ii) is the remark closing the
three-constituent part of \cref{sec:closure}, and the corresponding statement in \cref{sec:orth}. That $m$ and $\kappa$ are independent is
not a consequence of it: both are free inputs. What remains is an identification, and it is added, not derived. With $E=\hbar\times$rate,
$\hbar\kappa=mc^2$ holds if and only if $\ell=\hbar/(mc)$, the reduced Compton wavelength. The dynamics does not supply it. The interior equation
$\ddot X=-\kappa^2X$ does not contain $m$, which multiplies the Lagrangian; and the rescaling $(\tau,\kappa,I_0,m)\mapsto(s\tau,\kappa/s,sI_0,sm)$ of
\cref{cl:scalings}, which maps solutions to solutions and preserves the action, hence any quantisation of $h$, keeps $m\kappa$ and multiplies
$m/\kappa$ by $s^2$, so that $m=\hbar\kappa/c^2$ picks one point of each orbit by fiat. The rate is moreover ambiguous by a factor two: $q$ turns at
$2\kappa$ while $z$ turns at $\kappa$, and matching $2\hbar\kappa=mc^2$ instead gives $\ell=2\hbar/(mc)$. The algebra decides neither.

\subsection{Scale freedom}
\label{sec:restscale}

The Clifford data $\eta$, $G$, $G\eta^{-1}$, $\eta G$, the operator $\cK$ of \cref{lem:Gclifford} and the structure constants of the complex cross
product $\beta_0$ of \cref{prop:lGhoms} are integer matrices, and the rescaling $(I_0,m)\mapsto(sI_0,tm)$, $s,t>0$, preserves admissibility and
multiplies $\varepsilon$ by $1/(st)$ at fixed initial rate $\kappa_0$, by $s/t$ at fixed $h(0)$ (the remark following \cref{prop:several}). Two
further scalings act on the time.

\begin{proposition}[Scalings]
\label{cl:scalings}
\begin{enumerate}[label=(\roman*),leftmargin=2.2em]
\item The interior law $\ddot X=-\kappa^2X$ is invariant under $X\mapsto\sigma X$, which multiplies $\tfrac m2\eta(\alpha_+X,\alpha_+X)$ by
$\sigma^2$, absorbed by $m\mapsto m/\sigma^2$; and under $(\tau,\kappa)\mapsto(s\tau,\kappa/s)$, which together with $m\mapsto sm$ preserves the
action.
\item The closure is invariant under $(\tau,I_0,m)\mapsto(s\tau,sI_0,sm)$, with $(\theta,h,p)$ unchanged as functions of the rescaled time;
$\kappa\mapsto\kappa/s$, while $m\kappa$, $\varepsilon$ and the action are invariant.
\end{enumerate}
\end{proposition}

\begin{proof}
(i) Linearity; $Y(\tau')=X(\tau'/s)$ solves $Y''=-(\kappa/s)^2Y$, the Lagrangian with $sm$ in place of $m$ is multiplied by $s^{-1}$, and
$d\tau'=s\,d\tau$. (ii) With $d/d\tau'=s^{-1}d/d\tau$, $I_0'=sI_0$ and $N'=N/s$, both equations \eqref{eq:LProtor} are reproduced, $H'=H/s$, and
$\kappa_0'=\kappa_0/s$ gives $\varepsilon'=\varepsilon$; $h$ is unchanged.
\end{proof}

\begin{proposition}[Nothing in the algebra fixes $\ell$]
\label{cl:noscale}
\begin{enumerate}[label=(\roman*),leftmargin=2.2em]
\item Schur's lemma \cite{Schur} fixes only the form $c\,\id$ of the conversion operator (\cref{thm:onec}), and $c$ converts units of time into
units of length.
\item No monomial $c^{n_1}\hbar^{n_2}$ is a length, a mass or a time; with a mass $m$ adjoined, the only length is $\hbar/(mc)$, up to a number.
\item The interior systems with different $\kappa>0$ are carried into each other by a rescaling of $\tau$ that fixes every Clifford datum
(\cref{cl:scalings}(i)), so no condition formulated in the Clifford data distinguishes one value of $\kappa$, or of $\ell=c/\kappa$; in the
closure none fixes $\varepsilon$.
\end{enumerate}
\end{proposition}

\begin{proof}
(i) The commutant of the vector representation of $\so3$ in $M_3(\R)$ is $\R\,\id$. (ii) In the basis $(\mathsf M,\mathsf L,\mathsf T)$ the
dimensions $[c]=(0,1,-1)$ and $[\hbar]=(1,2,-1)$ span a plane containing none of $(1,0,0)$, $(0,1,0)$, $(0,0,1)$; adjoining $[m]=(1,0,0)$ gives
an invertible matrix, with $(0,1,0)=-[c]+[\hbar]-[m]$. (iii) The data are integer matrices, the rescaling acts on $\tau$ and $\kappa$ only, and
$\varepsilon$ is moved by the rescaling of $(I_0,m)$.
\end{proof}

Items (i) and (ii) are inherited dimensional analysis, already implicit in \eqref{eq:kappa}, where the sum of $\partial_\lambda$ and $G$ requires a
scale; (iii) is fixed by the algebra. This is the one exact sense available here of the phrase ``$m$ breaks the symmetry of $c^2$'': $c$ carries no
scale, the shift $\kappa G$ needs a length $\ell$ that the algebra cannot supply, and through $\ell=\hbar/(mc)$ that length would be a mass. What
it would break is the change of the unit of $\tau$ (\cref{cl:scalings}), not a symmetry of a vacuum; that is a reading.

\subsection{The frame drifts}
\label{sec:restdrift}

\Cref{ex:drift} computes the drift of the frame in closed form, $\Delta x=(8\mathrm K(-2)-4\mathrm E(-2))e_1$ per full turn of the rotor. After half a
turn $h$ has returned and $p=-p_0$, since $e^{\pi G}=-\id$; after a full turn both have returned. By \cref{thm:freedrift} every rotating solution
with a nonzero rest energy of the momentum field has an unbounded frame, and by \cref{cl:nogen}(ii) the rest energy of the frame field grows with
it, as $\tfrac12k^2|\Delta x|^2$ after $k$ turns in the example. A constant nonzero rate and a bounded frame coexist, for a free rotor with
Lorentz-class inertia, exactly on null data (\cref{sec:closure}); at the equilibria of \cref{cl:flat} the rate is zero and the frame moves
uniformly; and with several constituents they coexist on the balance surface of \cref{prop:several}. In motion, the differentiation of the energy
unbinds the frame.

\subsection{A quantum rotor, conditionally}
\label{sec:restquantum}

\begin{proposition}[The half turn]
\label{cl:halfturn}
\begin{enumerate}[label=(\roman*),leftmargin=2.2em]
\item $e^{\pi G}=-\id$, $e^{2\pi G}=\id$, and for $p_0\neq0$ the leaf circle $\{e^{\theta G}p_0\}$ has period exactly $2\pi$.
\item For every symmetric $N$, $U(\theta+\pi)=U(\theta)$.
\item $\Gamma=e_1e_2\cdots e_6$ anticommutes with every $e_a$ and $\Gamma^2=\id$, so $\Ad_\Gamma=-\id$ on $V$ and $\Gamma$ implements $e^{\pi G}$ on
spinors. By \cref{cl:norel}(ii) the multiples of the half turn are the only elements of the rotor group that lift to $\mathrm{Pin}(3,3)$.
\end{enumerate}
\end{proposition}

\begin{proof}
(i) $(e^{\theta G}-\id)p_0=(\cos\theta-1)p_0+\sin\theta\,Gp_0$ with $p_0\perp Gp_0$ and $|Gp_0|=|p_0|$. (ii) $p(\theta+\pi)=-p(\theta)$ and $U$ is
quadratic. (iii) A product of all six anticommuting generators anticommutes with each of them, and
$(e_1\cdots e_6)^2=(-1)^{15}\prod_a\eta_{aa}=\id$.
\end{proof}

The proposition restates the remark on the lift of the rotor at the end of the three-constituent part of \cref{sec:closure} and makes its
ingredients explicit; it is fixed. The quantum statement it permits is conditional on an assumption.

\begin{assumption}[A quantised leaf]
\label{cl:quant}
The leaf $T^*S^1$ of \cref{cl:frozen}(ii) is quantised by wave functions with $\Phi(\theta+2\pi)=e^{2\pi i\delta}\Phi(\theta)$, $h=-i\hbar\,\partial_\theta$,
and the Casimirs of the leaf are treated as numbers.
\end{assumption}

\begin{proposition}[A conditional rate ladder]
\label{cl:ladder}
Under \cref{cl:quant}:
\begin{enumerate}[label=(\roman*),leftmargin=2.2em]
\item $h$ has spectrum $\hbar(\Z+\delta)$. On the leaves on which $U$ is constant ($R_0=0$, or duality-class inertia) the levels are
$E^{\rm rot}_n=\hbar^2(n+\delta)^2/(2I_0)$, plus the constant $\tfrac12p_0^{\!\top}Np_0$ in the duality class, the rates are
$\kappa_n=(n+\delta)\hbar/I_0$, and $E^{\rm rot}_{n+1}-E^{\rm rot}_n=\tfrac\hbar2(\kappa_n+\kappa_{n+1})$.
\item If $R_0>0$ the Hamiltonian is a Mathieu operator \cite{Mathieu,DLMF} whose potential couples $e^{i(n+\delta)\theta}$ only to
$e^{i(n\pm2+\delta)\theta}$, so for every $\delta$ the parity of $n$ is conserved; it labels the eigenvalues $\pm e^{i\pi\delta}$ of the half-turn
translation $\theta\mapsto\theta+\pi$. For $q_{\rm M}:=I_0R_0/(2m\hbar^2)\gg1$ the lowest gap in either parity sector is
$\hbar\Omega-\hbar^2/(2I_0)+O(q_{\rm M}^{-1/2}\hbar^2/I_0)$, and the two sectors are degenerate up to tunnelling.
\item If moreover $\delta=0$, and, as a separate assumption, the first level $h=\hbar$ is matched by $\hbar\kappa_1=mc^2$ with $\ell=\hbar/(mc)$, then
$I_0=\hbar^2/(mc^2)$, $M_{\rm rot}=m$ and $E^{\rm rot}_1=\tfrac12mc^2$.
\end{enumerate}
\end{proposition}

\begin{proof}
(i) Spectrum of $-i\hbar\partial_\theta$ on twisted periodic functions; in the duality class $U$ is the constant found above. (ii)
$\cos(2\theta+\varphi_0)$ shifts $n$ by $\pm2$, and $e^{i(n+\delta)(\theta+\pi)}=(-1)^ne^{i\pi\delta}e^{i(n+\delta)\theta}$. Dividing the Schr\"odinger
equation by $-\hbar^2/(2I_0)$ gives $\Phi''+\bigl(a_{\rm M}-2q_{\rm M}\cos(2\theta+\varphi_0)\bigr)\Phi=0$ with $a_{\rm M}=2I_0E/\hbar^2$, and
$a_{\rm M}\sim-2q_{\rm M}+2(2r+1)\sqrt{q_{\rm M}}-\tfrac18\bigl((2r+1)^2+1\bigr)$ as $q_{\rm M}\to\infty$ \cite[\S28.8]{DLMF}. The levels $r=0,1$ lie in
the same sector and differ by $4\sqrt{q_{\rm M}}-1$, while $\hbar\Omega=(\hbar^2/2I_0)\,4\sqrt{q_{\rm M}}$. (iii) $\hbar/I_0=\kappa_1=c/\ell=mc^2/\hbar$,
$M_{\rm rot}=I_0/\ell^2$ and $E^{\rm rot}_1=\hbar^2/(2I_0)$.
\end{proof}

Numerically, at $R_0/m=2000$ and $\hbar=I_0=1$ the lowest gap in each sector is $62.7415$ for $\delta=0$ and for $\delta=0.3$, against
$\hbar\Omega-\hbar^2/(2I_0)=62.7456$, and $62.7416$ with the next term of the expansion. The proposition is constructed. The quantisation of the
leaf exists: it is the one the orbit method assigns \cite{Kirillov04}, with $\delta$ a character of the stabiliser $2\pi\Z$ in the universal
cover of the rotor group. Nothing selects $\delta$, $I_0$ or the matchings, and whether the ladder is physical is open. Missing are a reason for
the assumption, that is a quantum theory of the momentum field, which the closure does not give and the split-form frame obstructs: quantised
mode by mode with positive norm, the frame is a sum of oscillators of energies $\pm\hbar\kappa(j+\tfrac12)$, with no ground state. Missing too are
a value of $I_0$; a coupling to matter, since only the $\Z_2$ of the half turn reaches spinors, as $\Gamma$, and nothing ties the parity of $n$ to
chirality; and any derivation of the two matchings in (iii). The level $E^{\rm rot}_1=\tfrac12mc^2$ is $K_{\rm rot}=\tfrac12h\kappa$ at $h=\hbar$: the
factor $\tfrac12$ is built in, and $mc^2$ appears only as $\hbar$ times the rate put in, neither as a level nor as a transition energy.
